\partie{DEUXIÈME PARTIE}{DÉROULEMENT DU STAGE}

% Contenu exigé par le cours (module 2, §8, second temps) : poste occupé, position dans
% l'organigramme, faits importants, objectifs et enjeux, méthodes et moyens, collaborateurs,
% responsabilités, résultats, difficultés et solutions. Éviter le récit au jour le jour.

\section{Accueil, poste occupé et cadre de la mission}

\subsection{Accueil et intégration}
À notre arrivée, nous avons été reçu par \InfoTuteurPro, \InfoTuteurProPoste, qui nous a
présenté l'entreprise, son règlement intérieur et les consignes de sécurité. Après une visite
du siège et du magasin, nous avons été conduit sur le chantier \InfoChantier{} et présenté à
l'équipe plomberie. Nous avons reçu nos équipements de protection individuelle (casque,
chaussures de sécurité, gants, gilet haute visibilité).

\subsection{Poste occupé et position dans l'organigramme}
Nous avons occupé un poste d'\textbf{aide-technicien plombier} au sein de la division plomberie
(figure~\ref{fig:organigramme}). Nous travaillions sous l'autorité directe du chef d'équipe
plomberie, avec deux à trois techniciens et des manœuvres. Le chef de division plomberie et le
conducteur des travaux contrôlaient régulièrement l'avancement et la qualité des ouvrages.

\subsection{Objectifs et enjeux de la mission}
La mission confiée à l'équipe plomberie était de réaliser, dans les délais du planning général,
les réseaux de distribution d'eau froide et d'eau chaude sanitaire des étages, puis de poser
les organes de régulation et de contrôle (réducteurs de pression, manomètres, vannes
d'isolement). Nos objectifs personnels étaient les suivants :
\begin{itemize}
  \item maîtriser la lecture des plans d'exécution et des schémas de principe ;
  \item maîtriser l'assemblage des tubes PPR par polyfusion ;
  \item comprendre le rôle et le réglage des réducteurs de pression ;
  \item participer aux essais d'étanchéité et à la mise en eau ;
  \item découvrir l'organisation d'un chantier de grande hauteur et la coordination entre lots.
\end{itemize}
L'enjeu pour l'entreprise était double : respecter le planning contractuel, dont dépendaient les
corps d'état suivants (carreleurs, peintres, menuisiers), et livrer des réseaux étanches du
premier coup, car toute reprise après la pose des revêtements est coûteuse.

\section{Méthodes et moyens}

\subsection{Organisation du travail}
Chaque semaine, le chef d'équipe recevait du conducteur des travaux les zones à traiter et les
quantités de matériel correspondantes. Le travail était organisé \textbf{étage par étage} :
traçage et repérage, perçage et pose des supports, préparation des tubes, assemblage,
raccordement, puis essai. Une fiche d'autocontrôle était remplie à la fin de chaque zone.

\subsection{Moyens humains et matériels}
\begin{table}[H]\centering
\caption{moyens humains et matériels mobilisés}\label{tab:moyens}
\begin{tabularx}{\textwidth}{>{\bfseries}p{3.4cm} Y}
\toprule
\entete Catégorie & Détail \\
\midrule
Personnel & 1 chef d'équipe, 2 à 3 techniciens plombiers, 2 manœuvres, 1 stagiaire \\
Outillage & Polyfuseur et matrices de 20 à 63 mm, coupe-tubes, clés à molette et clés plates, niveau, mètre, perceuse-perforateur, pompe d'épreuve manuelle \\
Matériaux & Tubes et raccords PPR, flexibles inox, vannes à sphère, réducteurs de pression, manomètres, colliers et supports, isolant à cellules fermées \\
Sécurité & Casque, gants, chaussures de sécurité, lunettes, harnais pour les travaux en hauteur \\
Documents & Plans d'exécution, schémas de principe, fiches techniques des fabricants, fiches d'autocontrôle \\
\bottomrule
\end{tabularx}
\end{table}

\subsection{Personnes avec lesquelles nous avons collaboré}
Nous avons travaillé au quotidien avec le chef d'équipe plomberie et les techniciens de la
division. Nous avons également échangé avec l'équipe climatisation, notamment pour la
coordination des passages de tuyauteries dans les gaines techniques, et avec le magasinier
pour la préparation des sorties de matériel.

\section{Planning des activités menées}

Le tableau~\ref{tab:planning} présente le planning des activités, regroupées par période. La
figure~\ref{fig:gantt} en donne une vue d'ensemble.

\begin{longtable}{p{1.7cm} p{6.0cm} p{4.0cm} p{2.3cm}}
\caption{planning des activités menées}\label{tab:planning}\\
\toprule
\entete \textbf{Semaines} & \textbf{Activités} & \textbf{Outils} & \textbf{Domaine}\\
\midrule\endfirsthead
\toprule
\entete \textbf{Semaines} & \textbf{Activités} & \textbf{Outils} & \textbf{Domaine}\\
\midrule\endhead
\bottomrule\endfoot
S1 & Accueil, sécurité, visite du chantier, lecture des plans & Plans, EPI & Général \\
S2 -- S3 & Inventaire et préparation du matériel, traçage des réseaux aux étages & Mètre, niveau, cordeau & Plomberie \\
S4 -- S6 & Pose des supports et des canalisations PPR d'eau froide et d'eau chaude aux étages supérieurs & Polyfuseur, coupe-tube, perforateur & Plomberie \\
S7 -- S8 & Pose des canalisations aux étages inférieurs, raccordement aux colonnes montantes & Polyfuseur, coupe-tube & Plomberie \\
S9 & Essais d'étanchéité par zone, reprise des fuites & Pompe d'épreuve, manomètre & Plomberie \\
S10 -- S11 & Isolation des canalisations d'eau chaude et des réseaux exposés & Cutter, colle néoprène & Plomberie / Froid \\
S12 -- S13 & Montage des réducteurs de pression et des manomètres dans les locaux techniques et aux étages & Clés à molette, ruban PTFE & Plomberie \\
S14 & Pose des flexibles et de la robinetterie dans les salles d'eau & Clés plates, flexibles & Plomberie \\
S15 & Appui à l'équipe climatisation : pose de liaisons frigorifiques et d'évacuations de condensats & Cintreuse, dudgeonnière & Froid \\
S16 & Mise en eau, réglage des pressions, autocontrôles, bilan du stage & Manomètre & Plomberie \\
\end{longtable}

\begin{figure}[H]\centering
% Diagramme de Gantt du stage (16 semaines), cohérent avec le tableau du planning
\begin{tikzpicture}[x=0.72cm,y=0.55cm,font=\scriptsize]
  \foreach \s in {1,...,16} {
    \node[anchor=south] at (\s-0.5,0.2) {S\s};
    \draw[gray!30] (\s,0) -- (\s,-10);
  }
  \draw[gray!30] (0,0) -- (0,-10);
  \foreach [count=\i] \lib/\deb/\fin/\col in {
      Accueil et sécurité/0/1/vertfonce,
      Préparation et traçage/1/3/vertfonce,
      Canalisations étages hauts/3/6/vertmoyen,
      Canalisations étages bas/6/8/vertmoyen,
      Essais d'étanchéité/8/9/acompleter,
      Isolation des réseaux/9/11/vertfonce,
      Réducteurs et manomètres/11/13/vertmoyen,
      Robinetterie et flexibles/13/14/vertmoyen,
      Appui climatisation/14/15/bandeau!60!blue,
      Mise en eau et bilan/15/16/acompleter} {
    \node[anchor=east] at (-0.15,-\i+0.5) {\lib};
    \fill[\col, rounded corners=1.5pt] (\deb+0.05,-\i+0.15) rectangle (\fin-0.05,-\i+0.85);
  }
\end{tikzpicture}

\caption{vue d'ensemble du planning du stage}\label{fig:gantt}
\end{figure}

\section{Responsabilités, faits marquants et résultats}

\subsection{Responsabilités exercées}
Après une période d'observation, nous avons progressivement pris en charge des tâches
complètes : coupe et assemblage des tubes PPR, montage des réducteurs de pression et des
manomètres sur les colonnes des étages, isolation des réseaux d'eau chaude. Nous étions chargé
de la tenue de la fiche d'autocontrôle de notre zone et de la vérification de la présence du
matériel avant le début des travaux.

\subsection{Faits marquants}
\begin{itemize}
  \item La \textbf{détection d'une fuite} sur un assemblage lors d'un essai de pression : elle
        provenait d'une soudure réalisée avec un temps de chauffe insuffisant. L'assemblage a été
        coupé et refait. Cet incident a montré l'importance du respect des paramètres de soudage.
  \item La constatation d'une \textbf{pression statique trop élevée} aux étages inférieurs, qui
        justifie la pose de réducteurs de pression (voir troisième partie).
  \item Une \textbf{rupture d'approvisionnement} en raccords de 40 mm, qui a retardé une zone de
        deux jours et a conduit l'équipe à réorganiser l'ordre des étages.
\end{itemize}

\subsection{Résultats obtenus}
À la fin du stage, notre équipe avait achevé la distribution d'eau froide et d'eau chaude de
\AC{nombre} niveaux, monté \AC{nombre} réducteurs de pression et manomètres, et isolé environ
\AC{longueur} mètres de canalisations. Tous les réseaux auxquels nous avons participé ont passé
avec succès l'essai d'étanchéité après reprise.

\section{Difficultés rencontrées et solutions apportées}

\begin{table}[H]\centering
\caption{difficultés rencontrées et solutions apportées}\label{tab:difficultes}
\begin{tabularx}{\textwidth}{Y Y}
\toprule
\entete Difficulté & Solution apportée \\
\midrule
Manque de maîtrise initiale de la polyfusion (soudures irrégulières) & Entraînement sur des chutes de tube, respect strict des temps de chauffe affichés par le fabricant \\
Lecture des plans d'exécution de plusieurs niveaux & Accompagnement du chef d'équipe, repérage systématique des cotes et des niveaux \\
Ruptures de stock de certains raccords & Anticipation des besoins par zone, réorganisation de l'ordre des étages \\
Travail en hauteur et zones mal éclairées & Port du harnais, demande de projecteurs de chantier \\
Distance et coût du transport vers le chantier & Organisation du trajet en covoiturage avec des collègues \\
\bottomrule
\end{tabularx}
\end{table}

\section{Compétences acquises}

\subsection{Compétences professionnelles}
\begin{itemize}
  \item Lire un plan d'exécution et un schéma de principe de plomberie ;
  \item Couper, chanfreiner et assembler des tubes PPR par polyfusion ;
  \item Monter et régler un réducteur de pression, installer un manomètre ;
  \item Réaliser un essai d'étanchéité et interpréter la chute de pression ;
  \item Isoler une canalisation avec un isolant à cellules fermées ;
  \item Appliquer les règles de sécurité d'un chantier de grande hauteur.
\end{itemize}

\subsection{Compétences personnelles}
\begin{itemize}
  \item Ponctualité et rigueur dans l'exécution des tâches ;
  \item Travail en équipe et communication avec les autres corps d'état ;
  \item Autonomie et prise d'initiative ;
  \item Gestion du temps et respect des consignes.
\end{itemize}

\conclusionPartielle
Ce stage nous a permis d'occuper une place réelle dans une équipe de production. En passant de
l'observation à l'exécution, nous avons acquis des gestes professionnels précis et compris les
contraintes d'un chantier : délais, coordination, qualité et sécurité. Les difficultés
rencontrées ont été autant d'occasions d'apprendre. La troisième partie approfondit le thème
technique au cœur de ces activités : la plomberie sanitaire.
